\documentclass[11pt,a4paper]{article}
\usepackage[utf8]{inputenc}
\usepackage[T1]{fontenc}
\usepackage[french]{babel}
\usepackage[top=2cm,bottom=2cm,left=2cm,right=2cm]{geometry}
\usepackage{xcolor}
\usepackage{tcolorbox}
\tcbuselibrary{skins,breakable}
\usepackage{fancyhdr}
\usepackage{enumitem}
\usepackage{lmodern}

% --- Couleur discipline : ÉCONOMIE ---
\definecolor{mainColor}{RGB}{0,80,160}
\definecolor{darkGray}{RGB}{50,50,50}

\tcbset{boxrule=0.5pt, arc=3pt, left=5pt, right=5pt, top=3pt, bottom=3pt, breakable}

\newtcolorbox{defbox}[1]{
  colback=mainColor!8, colframe=mainColor, coltitle=white,
  fonttitle=\bfseries\footnotesize, title=DÉFINITION \textemdash\ #1}

\newtcolorbox{keybox}{
  colback=darkGray!7, colframe=darkGray!50, coltitle=white,
  fonttitle=\bfseries\footnotesize, title=POINT CLÉ}

\newtcolorbox{exbox}{
  colback=mainColor!5, colframe=mainColor!40,
  fonttitle=\bfseries\footnotesize, title=EXEMPLE}

\pagestyle{fancy}\fancyhf{}
\fancyhead[L]{\small\textcolor{mainColor}{\textbf{CEJM — BTS SIO}}}
\fancyhead[C]{\small\textbf{Chapitre 2 — Le fonctionnement des marchés}}
\fancyhead[R]{\small\textcolor{mainColor}{\textbf{ÉCONOMIE}}}
\fancyfoot[C]{\small\thepage}
\renewcommand{\headrulewidth}{0.5pt}
\setlength{\headheight}{15pt}
\setlist{itemsep=2pt, topsep=3pt, parsep=0pt}

\begin{document}

\begin{center}
  {\Large\bfseries\textcolor{mainColor}{Chapitre 2}}\\[2pt]
  {\large Le fonctionnement des marchés}\\[4pt]
  \textit{\textcolor{darkGray}{Comment se forme le prix d'équilibre et comment le marché s'autorégule-t-il ?}}
\end{center}
\vspace{4pt}\hrule\vspace{8pt}

\section*{Définitions essentielles}

\begin{defbox}{Marché}
Lieu de rencontre (réel ou virtuel) entre l'\textbf{offre} (vendeurs) et la \textbf{demande} (acheteurs) permettant de déterminer un prix.
\end{defbox}

\vspace{4pt}
\begin{defbox}{Prix d'équilibre}
Prix auquel \textbf{offre = demande}. Il n'y a ni excédent ni pénurie sur le marché.
\end{defbox}

\vspace{4pt}
\begin{defbox}{Élasticité-prix de la demande}
Mesure la \textbf{sensibilité de la demande} aux variations de prix.
$e = \frac{\Delta Q / Q}{\Delta P / P}$. Si $|e| > 1$ : demande élastique.
\end{defbox}

\section*{Les lois fondamentales du marché}

\subsection*{Loi de l'offre}
\begin{keybox}
\textbf{Prix $\uparrow$ → Offre $\uparrow$} (relation positive) : les producteurs sont incités à produire davantage.
\end{keybox}

\subsection*{Loi de la demande}
\begin{keybox}
\textbf{Prix $\uparrow$ → Demande $\downarrow$} (relation négative) : les consommateurs achètent moins quand le prix augmente.
\end{keybox}

\subsection*{L'équilibre de marché}
\begin{keybox}
L'équilibre se forme à l'\textbf{intersection} des courbes d'offre et de demande → autorégulation par le mécanisme des prix.
\end{keybox}

\section*{Les déséquilibres}

\begin{center}
\begin{tabular}{|l|l|l|}
\hline
\textbf{Situation} & \textbf{Cause} & \textbf{Conséquence} \\
\hline
Excédent & Offre > Demande & Prix tend à \textbf{baisser} \\
Pénurie & Demande > Offre & Prix tend à \textbf{monter} \\
\hline
\end{tabular}
\end{center}

\section*{Facteurs de déplacement des courbes}

\subsection*{Facteurs influençant la demande}
\begin{itemize}
  \item \textbf{Revenus des consommateurs} : revenus $\uparrow$ → demande $\uparrow$
  \item \textbf{Prix des biens substituts} : prix substitut $\uparrow$ → demande $\uparrow$
  \item \textbf{Prix des biens complémentaires} : prix complémentaire $\uparrow$ → demande $\downarrow$
  \item \textbf{Goûts et anticipations} : effet de mode, publicité
\end{itemize}

\subsection*{Facteurs influençant l'offre}
\begin{itemize}
  \item \textbf{Coûts de production} : coûts $\downarrow$ → offre $\uparrow$
  \item \textbf{Progrès technique} : amélioration → offre $\uparrow$
  \item \textbf{Nombre de producteurs} : plus de concurrents → offre $\uparrow$
  \item \textbf{Réglementation} : contraintes → offre $\downarrow$
\end{itemize}

\section*{Types de structures de marché}

\begin{center}
\begin{tabular}{|l|l|l|}
\hline
\textbf{Structure} & \textbf{Vendeurs} & \textbf{Prix} \\
\hline
Concurrence parfaite & Nombreux (atomicité) & Prix fixé par le marché \\
Monopole & 1 seul & Prix fixé par le monopoleur \\
Oligopole & Quelques-uns & Interdépendance des prix \\
Conc. monopolistique & Nombreux, produits diff. & Prix influencé (marque) \\
\hline
\end{tabular}
\end{center}

\begin{keybox}
\textbf{5 conditions de la concurrence parfaite} : atomicité, homogénéité, transparence, libre entrée/sortie, mobilité des facteurs.
\end{keybox}

\vspace{4pt}
\begin{exbox}
\textbf{Marché de l'essence} : oligopole (quelques compagnies), prix similaires, demande peu élastique (indispensable).\\
\textbf{Marché agricole} : proche de la concurrence parfaite, prix très volatils selon la récolte.
\end{exbox}

\section*{Points clés à retenir}

\begin{keybox}
Le \textbf{prix} est le signal central qui guide les décisions des agents économiques.
\end{keybox}
\vspace{3pt}
\begin{keybox}
En concurrence parfaite, les déséquilibres sont \textbf{temporaires} : le marché s'autorégule automatiquement.
\end{keybox}
\vspace{3pt}
\begin{keybox}
Les marchés réels s'éloignent souvent du modèle théorique → intervention de l'État souvent nécessaire.
\end{keybox}

\section*{Mots-clés}
\textcolor{mainColor}{\textbf{Marché}} $\cdot$
\textcolor{mainColor}{\textbf{Offre / Demande}} $\cdot$
\textcolor{mainColor}{\textbf{Prix d'équilibre}} $\cdot$
\textcolor{mainColor}{\textbf{Excédent / Pénurie}} $\cdot$
\textcolor{mainColor}{\textbf{Élasticité-prix}} $\cdot$
\textcolor{mainColor}{\textbf{Concurrence parfaite}} $\cdot$
\textcolor{mainColor}{\textbf{Monopole / Oligopole}} $\cdot$
\textcolor{mainColor}{\textbf{Biens substituts / complémentaires}} $\cdot$
\textcolor{mainColor}{\textbf{Autorégulation}}

\end{document}
